\documentclass[a4paper,onecolumn,10pt]{article}

\usepackage[margin=2cm]{geometry}
\usepackage[T1]{fontenc}
\usepackage{amsmath}
\usepackage{amssymb}
\usepackage{booktabs}
\usepackage{multirow}
\usepackage{microtype}
\usepackage[hidelinks]{hyperref}

\newcommand{\HCT}{\mathrm{HCT}}
\newcommand{\Beam}{\mathrm{Beam}}
\newcommand{\CISD}{\mathrm{CISD}}
\newcommand{\FCI}{\mathrm{FCI}}
\newcommand{\MPS}{\mathrm{MPS}}
\newcommand{\MPO}{\mathrm{MPO}}
\newcommand{\Ha}{\mathrm{Ha}}

\title{Orbital Localization, Clifford Frames, and Fiedler Ordering\\
in DMRG Benchmarks for N$_2$ and Linear H$_8$}
\author{}
\date{September 2026}

\begin{document}
\maketitle

\begin{abstract}
We compare fermionic spin-adapted and qubit-space density matrix
renormalization group (DMRG) calculations for correlated N$_2$ and linear
H$_8$ in the STO-3G basis.  The qubit Hamiltonians are considered before and
after Clifford transformations determined by two approximate-symmetry
selection procedures: Hamiltonian coefficient thresholding (HCT) and Beam
search.  Four spatial-orbital bases---canonical, split Pipek--Mezey, split
Pipek--Mezey restricted by point-group irrep, and CISD natural
orbitals---are tested with their original ordering and with a Fiedler ordering
derived from CISD mutual information.  Clifford transformation substantially
reduces the qubit-MPS bond dimension needed for chemical accuracy.  For N$_2$,
the best transformed calculations require bond dimension 8, compared with 39
for the canonical raw-qubit Hamiltonian.  For linear H$_8$, split
Pipek--Mezey localization and Fiedler ordering act cooperatively: the required
bond dimensions are 8, 57, 30, and 19 for fermionic, raw-qubit, HCT, and Beam
calculations, respectively.  The specialized fermionic SU(2) implementation
remains faster than all qubit calculations.  The results also show that a
smaller state bond dimension need not imply a shorter calculation when the
Clifford transformation substantially increases the MPO bond dimension.
\end{abstract}

\section{Scope of the comparison}

For each molecular system and orbital basis, we benchmark four Hamiltonian
frames:
\begin{enumerate}
    \item a fermionic quantum-chemistry Hamiltonian treated with SU(2)-adapted
          DMRG;
    \item the untransformed Jordan--Wigner qubit Hamiltonian;
    \item the qubit Hamiltonian transformed by a Clifford circuit synthesized
          from HCT approximate symmetries; and
    \item the qubit Hamiltonian transformed by a Clifford circuit synthesized
          from Beam-search approximate symmetries.
\end{enumerate}
Every frame is tested in its standard site order and in a Fiedler order
obtained from the corresponding CISD state.  The primary outcome is the
smallest tested maximum MPS bond dimension that simultaneously reaches
chemical accuracy and satisfies the DMRG sweep-energy convergence criterion.
The fermionic and qubit bond dimensions characterize different tensor-network
representations and should not be interpreted as equal computational
resources: the fermionic calculation uses four-state spatial-orbital sites and
SU(2) block sparsity, whereas the qubit calculation uses two-state
spin-orbital sites without particle-number or spin adaptation.

\section{Computational methods}

\subsection{Molecules and electronic-structure calculations}

All calculations use the STO-3G basis and retain all spatial orbitals and
electrons.  N$_2$ is evaluated at an N--N distance of 1.4~\AA; it contains 14
electrons in 10 spatial orbitals, corresponding to 20 qubits.  H$_8$ is a
linear, equally spaced chain with adjacent H--H separation 2.0~\AA; it
contains 8 electrons in 8 spatial orbitals, corresponding to 16 qubits.  The
chains are centered on the Cartesian $z$ axis.  Restricted Hartree--Fock
(RHF) calculations are performed in the Abelian $D_{2h}$ subgroup with an SCF
energy tolerance of $10^{-12}$ and a maximum of 100 cycles.  Spin orbitals are
interleaved as $(\alpha_0,\beta_0,\alpha_1,\beta_1,\ldots)$ before the
Jordan--Wigner transformation.

For every orbital basis, the one- and two-electron integrals, CISD state, and
FCI reference energy are recomputed.  FCI invariance under the orbital
rotations is verified to $10^{-8}$~Ha, and the FCI state is required to satisfy
$\langle \hat S^2\rangle<10^{-7}$.  The canonical reference energies are
listed in Table~\ref{tab:electronic-structure}.  CISD is not invariant to the
occupied--virtual mixing used to define the natural orbitals; its energy is
therefore recomputed rather than copied from the canonical calculation.

\begin{table}[htbp]
\centering
\caption{Electronic-structure dimensions and canonical-orbital reference
energies.  All energies are in Hartree.}
\label{tab:electronic-structure}
\begin{tabular}{lrrrrrr}
\toprule
System & $N_e$ & $N_{\rm orb}$ & $N_q$ & $E_{\rm RHF}$ & $E_{\rm CISD}$ & $E_{\rm FCI}$ \\
\midrule
N$_2$ & 14 & 10 & 20 & $-107.3578154453$ & $-107.5832737773$ & $-107.6231741781$ \\
H$_8$ &  8 &  8 & 16 & $  -3.1614329658$ & $  -3.5542953056$ & $  -3.7966934506$ \\
\bottomrule
\end{tabular}
\end{table}

\subsection{Orbital bases}

Four spatial-orbital bases are compared.
\begin{description}
    \item[Canonical:] symmetry-adapted RHF molecular orbitals.
    \item[Split PM:] Pipek--Mezey localization is applied separately to the
          canonical occupied and virtual subspaces.  The occupied--virtual
          partition is not mixed.
    \item[Split PM, irrep-restricted:] each occupied or virtual block is
          divided further according to its canonical $D_{2h}$ irrep, and
          Pipek--Mezey localization is performed independently in every
          nontrivial block.
    \item[Natural:] the spin-summed spatial one-particle density matrix of the
          canonical CISD state is diagonalized globally.  The resulting
          orbitals are sorted by decreasing occupation.  This transformation
          may mix canonical occupied and virtual orbitals; the CISD state is
          consequently recomputed in the natural-orbital basis.
\end{description}
The Pipek--Mezey convergence tolerance is $10^{-10}$ with at most 1000
localization cycles.  Orbital signs are fixed deterministically from the
largest-magnitude AO coefficient.  Orthonormality is verified to $10^{-9}$;
for the split localizations, preservation of the occupied projector is
verified to $10^{-8}$.

\subsection{Pauli Hamiltonians and approximate symmetries}

The fermionic Hamiltonian is mapped directly to a packed Pauli representation
with the repository's streamed Jordan--Wigner implementation.  Equal Pauli
words are combined, and combined coefficients with magnitude at most
$10^{-12}$ are discarded.  A Pauli word is stored as its binary $X$ and $Z$
masks, avoiding construction of an intermediate OpenFermion
\texttt{QubitOperator}.  The stream retains the first-occurrence order induced
by the fermionic-term order and its Jordan--Wigner expansion.  No subsequent
sorting by coefficient magnitude, Pauli weight, or spatial support is used.

For a Pauli product $S$, both symmetry searches use the CISD squared-commutator
cost
\begin{equation}
    c(S)=
    \langle\Psi_{\CISD}|[H,S]^\dagger[H,S]|\Psi_{\CISD}\rangle.
    \label{eq:symmetry-cost}
\end{equation}
Each search returns $N_q$ linearly independent, mutually commuting Pauli
generators.  HCT progressively removes Hamiltonian terms by increasing a
coefficient threshold over the distinct coefficient magnitudes above
$10^{-5}$.  At each threshold, the symplectic null space is reduced to a
maximal commuting subspace, with candidate generators ordered first by
Eq.~\eqref{eq:symmetry-cost}, then by Pauli weight, and finally by their packed
binary representation.  Previously selected generators are retained as the
threshold increases.

The Beam candidate pool contains the HCT generators, generators obtained from
at most 256 Hamiltonian terms, and pairwise products seeded by the 50 leading
Hamiltonian terms.  The pool is sorted by increasing singleton cost and capped
at 1000 candidates; no maximum Pauli weight is imposed.  Beam search has width
16 and selects a full-rank commuting set by maximizing the negative sum of the
singleton costs, equivalently minimizing $\sum_S c(S)$.  The configured
heavy-core fraction is 0.95.  The resulting total costs are reported in
Table~\ref{tab:symmetry-costs}; Fiedler reordering does not alter the selected
symmetry set or its cost.

\begin{table}[htbp]
\centering
\caption{Total CISD squared-commutator cost $\sum_S c(S)$ of the selected
full-rank symmetry sets.}
\label{tab:symmetry-costs}
\begin{tabular}{llrr}
\toprule
System & Orbital basis & HCT & Beam \\
\midrule
\multirow{4}{*}{N$_2$}
 & Canonical                  & 2.724 & 2.090 \\
 & Split PM                   & 3.602 & 2.779 \\
 & Split PM, irrep-restricted & 2.691 & 2.091 \\
 & Natural                    & 2.643 & 2.070 \\
\midrule
\multirow{4}{*}{H$_8$}
 & Canonical                  & 6.224 & 3.409 \\
 & Split PM                   & 1.253 & 0.617 \\
 & Split PM, irrep-restricted & 4.425 & 2.987 \\
 & Natural                    & 5.968 & 3.514 \\
\bottomrule
\end{tabular}
\end{table}

For each selected set, a Clifford circuit maps the generators to positive
single-qubit $Z$ operators.  The same circuit is applied to the Hamiltonian and
the CISD state, so the spectrum and the warm-start expectation value are
preserved.  Clifford action on the sparse state is evaluated exactly by
bit-string and coefficient updates before conversion to an MPS; no determinant
truncation is applied during the transformation.

\subsection{Fiedler ordering}

Fiedler orderings are derived from the CISD state separately for every
Hamiltonian frame.  For raw-qubit calculations, the graph vertices are
individual spin-orbital qubits.  For the fermionic calculation, each vertex is
a four-state spatial orbital containing its paired $\alpha$ and $\beta$ modes.
For HCT and Beam frames, the CISD state is first transformed by the
corresponding Clifford circuit and the mutual information is then evaluated in
that frame.  One- and two-site reduced density matrices are contracted from the
sparse state, entropies are evaluated in base 2, and the standard mutual
information convention is used.  Edges below $10^{-12}$ are omitted.  Sites
are ordered by the Fiedler vector of the weighted graph Laplacian, with index
tie-breaking and disconnected components ordered by total edge weight.  The
resulting permutation is applied identically to the Hamiltonian, symmetry
operators, and CISD state.

\subsection{MPO and MPS construction}

Fermionic MPOs are built directly from the spatial one- and two-electron
integrals with the SU(2)-adapted quantum-chemistry MPO builder in
\texttt{pyblock2}; no nontrivial point-group labels are supplied to Block2.
Qubit MPOs are built from the Pauli stream with Block2's
\texttt{FastBlockedSumBipartite} algorithm, using a block parameter of 10 and
an MPO compression cutoff of $10^{-10}$, without an explicit MPO bond cap.
Pauli products are supplied in their deterministic stream order with symbolic
reordering disabled.

The sparse CISD state is converted to an MPS by adding determinants in
decreasing coefficient magnitude in batches of 32 and compressing with an SVD
cutoff of $10^{-13}$.  The temporary bond cap during construction and Clifford
transformation equals the largest bond dimension in the corresponding DMRG
search.  The constructed MPS is normalized, and its expectation value with
the final MPO must agree with the independently computed CISD energy within
$10^{-7}$~Ha when the accumulated construction loss is below $10^{-12}$.

\subsection{DMRG convergence search}

All fermionic and qubit-space DMRG calculations use two-site updates in
\texttt{pyblock2}.  Each tested maximum bond dimension $\chi$ starts from a
fresh copy of the same CISD warm-start MPS rather than from the result at a
neighboring bond dimension.  DMRG is run for at most 100 sweeps with fixed
maximum bond dimension $\chi$.  The noise schedule is
\begin{equation}
    (10^{-4},10^{-4},10^{-5},10^{-5},10^{-6},10^{-6},0,\ldots).
\end{equation}
The local Davidson threshold is $10^{-10}$, with at most 50 Davidson
iterations per local optimization.  Sweep convergence requires that the
absolute change between the final two sweep energies be below $10^{-6}$~Ha
after the noise has been removed.  A bond dimension is accepted only if this
condition and chemical accuracy both hold:
\begin{equation}
    |E_{\rm DMRG}(\chi)-E_{\FCI}|
    \leq 1.6\times10^{-3}\ \Ha.
    \label{eq:chemical-accuracy}
\end{equation}

An integer bisection search locates the smallest accepted $\chi$ in
$1\leq\chi\leq200$ for N$_2$ and $1\leq\chi\leq300$ for H$_8$.  The N$_2$
run used the midpoint-first implementation; the completed H$_8$ data used the
earlier endpoint-bracketed implementation.  In every successful frame the
immediately lower integer was evaluated and failed, and no evaluated frame
showed a successful smaller bond dimension followed by a failed larger one.
This establishes the reported minimum within the individual searches, subject
to the usual possibility of DMRG convergence to different local stationary
states.

The runtime reported in the result tables is the wall time of the DMRG
optimization at the accepted bond dimension, summed over all completed sweeps.
It excludes electronic-structure calculations, symmetry search, Fiedler
ordering, MPO construction, initial-MPS construction, and DMRG calculations
at other bond dimensions.  Per-sweep and complete search timings are retained
in the saved JSON and CSV records.

\subsection{Software and execution environment}

The calculations used Python 3.13.5, NumPy 2.2.5, SciPy 1.15.3, PySCF 2.14.0,
OpenFermion 1.7.1, OpenFermion-PySCF 0.5, and Block2/pyblock2 0.5.3.  DMRG
used four Block2 threads and one BLAS/LAPACK thread.  The reported local timing
data were obtained on an Apple M1 Pro MacBook Pro with eight CPU cores and
16~GB memory; timings should therefore be interpreted as within-machine
comparisons.  The benchmark does not explicitly set Block2's internal random
seed.  The CISD warm start and all search tie-breaking are deterministic, but
exact runtime and convergence near the chemical-accuracy boundary can still
depend on the software build, threading, and internal random-number state.

\section{Results}

Tables~\ref{tab:n2-orbital-dmrg} and~\ref{tab:h8-orbital-dmrg} report the
smallest accepted MPS bond dimension and the corresponding DMRG optimization
time.  Boldface marks the smallest bond dimension within each Hamiltonian
representation for a given molecule; ties are all shown in bold.

\begin{table}[htbp]
\centering
\caption{N$_2$ at $R=1.4$~\AA.  Each entry is
$\chi_{\min}$ [DMRG wall time in seconds].}
\label{tab:n2-orbital-dmrg}
\begin{tabular}{llcccc}
\toprule
Orbital basis & Ordering & Fermionic SU(2) & Raw qubit & HCT & Beam \\
\midrule
\multirow{2}{*}{Canonical}
 & Standard & 23 [0.21] & 39 [11.49] & 12 [1.21] & 10 [2.69] \\
 & Fiedler  & 31 [0.42] & 32 [16.37] & 13 [1.89] & \textbf{8 [1.81]} \\
\midrule
\multirow{2}{*}{Split PM}
 & Standard & 23 [0.36] & 50 [13.19] & 27 [19.52] & 22 [169.73] \\
 & Fiedler  & \textbf{12 [0.47]} & 33 [17.46] & 21 [8.03] & 19 [40.65] \\
\midrule
\multirow{2}{*}{Split PM, irrep-restricted}
 & Standard & 23 [0.23] & 38 [14.60] & 10 [1.78] & 10 [1.76] \\
 & Fiedler  & 33 [0.37] & 31 [5.90] & 12 [1.76] & \textbf{8 [1.85]} \\
\midrule
\multirow{2}{*}{Natural}
 & Standard & 24 [0.22] & 35 [7.33] & \textbf{9 [1.33]} & 10 [1.91] \\
 & Fiedler  & 33 [0.37] & \textbf{28 [8.32]} & 10 [1.92] & 11 [1.51] \\
\bottomrule
\end{tabular}
\end{table}

\begin{table}[htbp]
\centering
\caption{Linear H$_8$ at adjacent H--H separation 2.0~\AA.  Each entry is
$\chi_{\min}$ [DMRG wall time in seconds].  A dagger marks a run that reached
the sweep-energy tolerance but not chemical accuracy by $\chi=300$; its
stationary higher-energy solution is consistent with optimization trapping
and is not interpreted as a bond-dimension lower bound.}
\label{tab:h8-orbital-dmrg}
\begin{tabular}{llcccc}
\toprule
Orbital basis & Ordering & Fermionic SU(2) & Raw qubit & HCT & Beam \\
\midrule
\multirow{2}{*}{Canonical}
 & Standard & 35 [0.37] & 123 [49.66] & ---$^\dagger$ & 47 [9.45] \\
 & Fiedler  & 35 [0.40] & 122 [49.05] & 54 [11.56] & 49 [10.71] \\
\midrule
\multirow{2}{*}{Split PM}
 & Standard & 32 [0.64] & 117 [54.96] & 37 [12.16] & 32 [16.73] \\
 & Fiedler  & \textbf{8 [0.23]} & \textbf{57 [18.68]}
             & \textbf{30 [8.63]} & \textbf{19 [6.34]} \\
\midrule
\multirow{2}{*}{Split PM, irrep-restricted}
 & Standard & 35 [0.38] & 123 [45.39] & ---$^\dagger$ & 45 [8.33] \\
 & Fiedler  & 33 [0.39] & 105 [38.36] & ---$^\dagger$ & 46 [9.50] \\
\midrule
\multirow{2}{*}{Natural}
 & Standard & 35 [0.34] & 123 [48.96] & 52 [11.35] & ---$^\dagger$ \\
 & Fiedler  & 35 [0.44] & 121 [85.95] & 51 [10.19] & 48 [7.32] \\
\bottomrule
\end{tabular}
\end{table}

\subsection{\texorpdfstring{N$_2$}{N2}: Clifford compression without a localization advantage}

For N$_2$, both Clifford constructions substantially reduce the qubit-MPS
bond dimension.  In canonical orbitals, the raw-qubit threshold is 39,
compared with 12 for HCT and 10 for Beam.  Beam with Fiedler ordering reaches
$\chi=8$, as does Beam in the irrep-restricted split-PM basis; the lowest HCT
threshold is $\chi=9$ in standard-order natural orbitals.  The CISD natural
orbitals also give the smallest raw-qubit threshold, 28, after Fiedler
ordering.

Unrestricted split-PM localization is unfavorable for the transformed N$_2$
calculations.  Relative to canonical orbitals, its HCT and Beam costs increase
by approximately 32--33\%, while the standard-order thresholds increase from
12 to 27 and from 10 to 22, respectively.  The corresponding Pauli MPOs also
grow sharply: the standard-order Beam MPO increases from bond dimension 832
to 6097.  Consequently, its accepted $\chi=22$ calculation requires
169.73~s, despite having a smaller state bond dimension than the raw-qubit
calculation.  Restricting localization to occupied/virtual irrep blocks
largely restores the canonical symmetry costs, MPO dimensions, and DMRG
thresholds.

Fiedler ordering is not uniformly beneficial for this compact diatomic.  It
improves the canonical raw-qubit and Beam thresholds and strongly improves the
split-PM fermionic result, but it increases several fermionic, HCT, or Beam
thresholds in the canonical, irrep-restricted, and natural bases.  Orbital
localization and ordering must therefore be evaluated together rather than
treated as independently favorable transformations.

\subsection{\texorpdfstring{H$_8$}{H8}: cooperative localization and ordering}

The extended H$_8$ chain shows the expected real-space localization benefit.
Split-PM localization lowers the HCT and Beam symmetry costs from 6.224 and
3.409 in canonical orbitals to 1.253 and 0.617.  Combining this basis with
Fiedler ordering gives the smallest threshold in every representation:
$\chi=8$ for fermionic SU(2), 57 for raw qubits, 30 for HCT, and 19 for Beam.
Thus Beam reduces the required qubit-MPS bond dimension by a factor of three
relative to the raw-qubit calculation in the same orbital basis and ordering.

The irrep-restricted localization yields much smaller improvements.  This
contrast indicates that mixing canonical point-group blocks is important for
forming orbitals localized along the hydrogen chain, whereas the same mixing
is detrimental to N$_2$.  Natural orbitals also provide little reduction for
raw H$_8$: the standard thresholds remain 123 in the canonical,
irrep-restricted, and natural bases.

Four transformed H$_8$ calculations satisfy the sweep-energy stopping
criterion but fail Eq.~\eqref{eq:chemical-accuracy}: standard-order canonical
HCT, both irrep-restricted HCT orderings, and standard-order natural Beam.  At
$\chi=300$, these calculations remain near $-3.6711575235$~Ha, approximately
125.5~mHa above FCI, and increasing the nominal bond cap does not lower the
energy.  Standard-order canonical HCT and natural Beam both reach chemical
accuracy after Fiedler reordering.  These observations identify local DMRG
optimization behavior, rather than an insufficient representational bond
dimension, as the likely origin of at least those two failures.

\subsection{Operator complexity and runtime}

Table~\ref{tab:mpo-dimensions} shows that Clifford compression of the state
can increase the MPO bond dimension.  The effect is particularly large for
unrestricted split-PM N$_2$, where the Beam MPO reaches 6097.  For split-PM
H$_8$ with Fiedler ordering, Beam reduces the state threshold from 57 to 19
while increasing the MPO bond dimension from 697 to 2200.  That tradeoff is
still favorable for the accepted DMRG calculation: the Beam calculation takes
6.34~s, compared with 18.68~s for raw qubits.  For N$_2$, however, the much
larger Beam MPO dominates the smaller state bond dimension and increases the
runtime.

\begin{table}[htbp]
\centering
\small
\caption{Maximum MPO bond dimensions.  Each entry gives standard/Fiedler
ordering.  The fermionic and qubit MPO bond dimensions refer to different
local site spaces and should be compared primarily within a representation.}
\label{tab:mpo-dimensions}
\begin{tabular}{llrrrr}
\toprule
System & Orbital basis & Fermionic SU(2) & Raw qubit & HCT & Beam \\
\midrule
\multirow{4}{*}{N$_2$}
 & Canonical                  & 63/87   & 291/575  & 539/861   & 832/730 \\
 & Split PM                   & 123/123 & 643/3752 & 3046/2561 & 6097/4178 \\
 & Split PM, irrep-restricted & 63/87   & 291/544  & 525/923   & 824/707 \\
 & Natural                    & 63/87   & 291/586  & 655/979   & 824/1150 \\
\midrule
\multirow{4}{*}{H$_8$}
 & Canonical                  & 83/83 & 349/690 & 2095/1855 & 1173/1226 \\
 & Split PM                   & 83/83 & 411/697 & 2069/1903 & 3355/2200 \\
 & Split PM, irrep-restricted & 83/83 & 349/846 & 1227/1532 & 1491/1382 \\
 & Natural                    & 83/83 & 349/693 & 1678/1496 & 1189/1154 \\
\bottomrule
\end{tabular}
\end{table}

The fermionic SU(2) implementation is the fastest method throughout the
successful runs.  At the accepted bond dimension, fermionic optimization
takes 0.21--0.47~s for N$_2$ and 0.23--0.64~s for H$_8$.  Successful Beam
calculations take 1.51--169.73~s for N$_2$ and 6.34--16.73~s for H$_8$.
This advantage reflects the specialized quantum-chemistry MPO, spatial-site
representation, and SU(2) block sparsity.  The Clifford benchmarks therefore
demonstrate reduced qubit-MPS entanglement requirements, not a wall-time
advantage over spin-adapted fermionic DMRG.

\section{Limitations and next steps}

The benchmark establishes trends for two STO-3G systems and a single CISD
warm start per frame.  The integer bisection assumes that DMRG acceptance is
monotone with the allowed bond dimension; the implementation detects
nonmonotone evaluated outcomes, but independent random seeds were not used to
quantify optimizer variability near the threshold.  The stationary H$_8$
failures show that sweep-energy convergence alone does not establish ground-
state convergence.  Future comparisons should repeat marginal and failed
frames with controlled seeds or alternative noise schedules.

The current Pauli terms are supplied to the MPO builder in streamed
Jordan--Wigner order.  Sorting by coefficient magnitude, spatial support, or
an MPO-specific graph heuristic may reduce intermediate and final MPO bond
dimensions without changing the Hamiltonian and is a direct target for
follow-up tests.  Additional directions are orbital optimization against an
entanglement- or MPO-aware objective, systematic Beam-pool and beam-width
studies, iterative searches restricted to Pauli-$Z$ quasi-symmetries, and
fermionic Gaussian transformations defined directly in Majorana space.

\section{Reproduction and saved data}

The authoritative driver is
\texttt{scripts/benchmark\_sto3g\_orbital\_dmrg.py}.  The following commands
specify the two calculations; all remaining parameters are the defaults
reported above.

\begin{verbatim}
python -u scripts/benchmark_sto3g_orbital_dmrg.py \
  --system N2_corr --orderings standard fiedler \
  --min-bond-dim 1 --max-bond-dim 200

python -u scripts/benchmark_sto3g_orbital_dmrg.py \
  --system H8_corr --orderings standard fiedler \
  --min-bond-dim 1 --max-bond-dim 300
\end{verbatim}

The complete result directories are
\begin{verbatim}
saved/results/n2_corr_sto3g_orbital_dmrg/
saved/results/h8_corr_sto3g_orbital_dmrg/
\end{verbatim}
Each directory contains the exact configuration in \texttt{settings.json},
aggregate thresholds in \texttt{dmrg\_summary.csv}, per-bond energies and
timings in \texttt{dmrg\_curves.csv} and JSON form, and one subdirectory per
orbital basis.  The basis subdirectories retain molecular integrals and
orbitals, sparse CISD states, selected Pauli generators and their individual
costs, Clifford/Fiedler metadata, MPO bond dimensions, per-sweep energies and
runtimes, and native Block2 MPO/MPS restart artifacts.  Results are
checkpointed after every completed bond-dimension calculation, and
\texttt{--resume} continues an interrupted calculation only when its
result-defining settings agree with the saved configuration.

\end{document}
